% =====================================================================
%  CLASE 11 -- Drenaje transversal: puentes
% =====================================================================
\clase{11}{Drenaje transversal de caminos y carreteras: puentes}

\section{Diseño hidráulico de puentes}
Cuando el caudal a atravesar es muy grande, la única alternativa es la
construcción de un puente.
\begin{itemize}
  \item \textbf{Opción 1: terraplén y estribos} (más económico). \emph{Método
        fusible}: el objetivo es que el río, en una crecida, se lleve el
        terraplén y no la estructura.
  \item \textbf{Opción 2: puente que cubre todo el cauce, con cepas} (¡alto costo!).
\end{itemize}

\begin{figure}[H]
\centering
\begin{subfigure}{0.48\textwidth}\centering
\begin{tikzpicture}[font=\small,scale=0.8]
  \fill[suelo!80] (0,3) .. controls (0.8,3) and (0.8,1.5) .. (2,1.3) -- (2.4,0.6) -- (4.4,0.6) -- (4.8,1.3) .. controls (6,1.5) and (6.2,3) .. (7,3) -- (7,0) -- (0,0) -- cycle;
  \fill[sueloclaro] (2,1.4) -- (2.4,0.6) -- (3.0,0.6) -- (3.0,1.4) -- cycle;
  \fill[sueloclaro] (3.9,1.4) -- (3.9,0.6) -- (4.4,0.6) -- (4.8,1.4) -- cycle;
  \fill[aguaclara] (3.0,0.6) rectangle (3.9,0.9);
  \draw[very thick,gray] (1.2,1.45) -- (5.8,1.45);
  \draw[thick] (3.0,0.6) -- (3.0,1.45) (3.9,0.6) -- (3.9,1.45);
  \node[azulusm] at (3.45,2.4) {Terraplén};
  \draw[->,azulusm] (3.1,2.2) -- (2.6,1.1); \draw[->,azulusm] (3.8,2.2) -- (4.3,1.1);
\end{tikzpicture}
\caption{Opción 1: terraplén y estribos.}
\end{subfigure}\hfill
\begin{subfigure}{0.48\textwidth}\centering
\begin{tikzpicture}[font=\small,scale=0.8]
  \fill[suelo!80] (0,3) .. controls (0.3,1.5) and (1.2,1.5) .. (1.8,0.6) -- (5.2,0.6) .. controls (5.8,1.5) and (6.7,1.5) .. (7,3) -- (7,0) -- (0,0) -- cycle;
  \fill[aguaclara] (3.1,0.6) rectangle (3.9,0.8);
  \fill[gray!60] (0,3) rectangle (7,3.2);
  \foreach \x in {1.2,5.8} \fill[gray!60] (\x-0.1,1.5) rectangle (\x+0.1,3);
\end{tikzpicture}
\caption{Opción 2: puente con cepas.}
\end{subfigure}
\caption{Alternativas para cruzar un cauce de gran caudal.}
\end{figure}

\subsection{Recomendaciones de diseño}
\begin{itemize}
  \item No generar resaltos.
  \item Evitar que el peraltamiento del escurrimiento aguas arriba produzca
        inundaciones que antes no ocurrían.
  \item Que no se generen velocidades excesivas que puedan provocar problemas
        de socavación.
\end{itemize}
(Ejemplo reciente: desborde del Canal San Carlos, 16 de junio de 2024.)

\section{Opción 1: terraplén y estribos}
Considerando escurrimiento de río en el cauce (ancho $b$) que se contrae al
ancho $L$ del puente. Se definen las secciones (1) aguas abajo, (2) en la
contracción y (3) aguas arriba.

\begin{figure}[H]
\centering
\begin{tikzpicture}[font=\small]
  \fill[top color=aguaclara,bottom color=agua] (0,0) rectangle (8,2.6);
  \fill[sueloclaro,draw] (3.2,1.9) rectangle (4.4,2.6);
  \fill[sueloclaro,draw] (3.2,0) rectangle (4.4,0.7);
  \fill[gray!60] (3.2,1.8) rectangle (4.4,1.9); \fill[gray!60] (3.2,0.7) rectangle (4.4,0.8);
  \draw[cota] (0.9,0) -- (0.9,2.6) node[midway,left]{$b$};
  \draw[cota] (3.8,0.8) -- (3.8,1.8) node[midway,left]{$L$};
  \foreach \x/\n in {2.6/3,4.4/2,5.4/1} {\draw[rojo,dashed,thick] (\x,-0.4) -- (\x,3.0); \node[rojo,below] at (\x,-0.4) {(\n)};}
  \node[azulusm] at (6.6,3.3) {Terraplén};
  \draw[->,azulusm] (6.2,3.1) -- (4.3,2.4); \draw[->,azulusm] (6.9,3.1) .. controls (7.8,-0.6) and (5,-0.6) .. (4.3,0.3);
  \draw[flecha,azulusm] (7.6,1.3) -- (6.9,1.3);
\end{tikzpicture}
\caption{Contracción del cauce por terraplenes y estribos (vista en planta;
el flujo va de (3) hacia (1)).}
\end{figure}

La momenta (fuerza específica) en cada sección es:
\begin{formula}
$\displaystyle M_1=\frac{Q^2}{g(b\cdot h_1)}+\frac{h_1}{2}(b\cdot h_1)\qquad
M_2=\frac{Q^2}{g(L\cdot h_2)}+\frac{h_2}{2}(b\cdot h_2)$
\end{formula}

\begin{figure}[H]
\centering
\begin{tikzpicture}
\begin{axis}[width=8cm,height=6cm,xmin=0,xmax=10,ymin=0,ymax=5,axis lines=left,
  xlabel={$M$},ylabel={$h$},xtick=\empty,ytick=\empty,clip=false,font=\small]
  \addplot[thick,azulusm,domain=0.35:4.6,samples=60] ({6/x+x^2/2*0.9},{x});
  \addplot[thick,rojo,domain=0.45:4.6,samples=60] ({9/x+x^2/2*0.9},{x});
  \addplot[thick,green!60!black,domain=0.5:4.6,samples=60] ({11/x+x^2/2*0.9},{x});
  \addplot[thick,orange,domain=0.55:4.6,samples=60] ({13/x+x^2/2*0.9},{x});
  \draw[dashed] (axis cs:0,3.4) node[left]{$h_{n1}$} -- (axis cs:{6/3.4+3.4^2*0.45},3.4) -- (axis cs:{6/3.4+3.4^2*0.45},0) node[below]{$M_1(h_{n1})$};
\end{axis}
\end{tikzpicture}
\caption{Curvas de momenta: sección 1 (azul) y sección 2 con $M_{c2}<M_1$
(rojo), $M_{c2}=M_1$ (verde) y $M_{c2}>M_1$ (naranja, se produce resalto).}
\end{figure}

Para la sección (2) se tiene:
\[
\frac{\partial M_2}{\partial h_2}=0\quad\Longrightarrow\quad h_{c2,M}=\sqrt[3]{\frac{Q^2}{g\cdot b\cdot L}}
\]
Por lo tanto, para definir el largo mínimo del puente ($L$), tal que no se
produzca resalto, se considera la condición crítica, resultando el siguiente
sistema de ecuaciones:
\begin{formula}
$M_1=M_{2c}\quad(1)\qquad\qquad\displaystyle h_{c2,M}=\sqrt[3]{\frac{Q^2}{g\cdot b\cdot L}}\quad(2)$
\end{formula}
donde
\[
M_{2c}=\frac{Q^2}{g(L\cdot h_{2c,M})}+\frac{h_{2c,M}}{2}(b\cdot h_{2c,M}),\qquad
M_1=\frac{Q^2}{g(b\cdot h_1)}+\frac{h_1}{2}(b\cdot h_1)\ \ \text{(conocido)}
\quad\Longrightarrow\quad\boxed{L\geq L_{min}}
\]

Con $L$ definido para que no se produzca resalto, se puede realizar un
balance de energía entre (1) y (2):
\begin{formula}
$\displaystyle h_2+\frac{v_2^2}{2g}=h_{n1}+\frac{v_1^2}{2g}+\Lambda_s$
\end{formula}
donde las pérdidas singulares a la salida se calculan a través de la
ecuación de Borda--Carnot:
\begin{formula}
$\displaystyle\Lambda_s=\frac{(v_2-v_1)^2}{2g}-\frac{(h_{n1}-h_2)^2}{2h_{n1}}$
\end{formula}
Además, $v_1=\dfrac{Q}{h_{n1}b}$ y $v_2=\dfrac{Q}{h_2L}$, de donde se obtiene $h_2$.

Luego, realizando un balance de energía entre (2) y (3):
\begin{formula}
$\displaystyle h_3+\frac{v_3^2}{2g}=h_2+\frac{v_2^2}{2g}+\Lambda_e,\qquad
v_3=\frac{Q}{h_3b},\qquad\Lambda_e=k_e\frac{v_2^2}{2g},\quad k_e=0.1\sim0.3$
\end{formula}
de donde se obtiene $h_3$. Se debe verificar que
\[
h_3<h_{max}\qquad\qquad v_2<v_{max}
\]
Si no se cumple alguna de las condiciones, se debe elegir un $L$ mayor y
recalcular todo hasta que se cumplan; o bien, a partir de dichas condiciones
se determinan $L$ y $h_2$ y se verifica que no haya resalto.

\paragraph{Conservación del momentum lineal entre (3) y (2).} También se
puede aplicar, con:
\[
n=\frac{b_3}{b_2},\qquad X_3=\frac{h_3}{h_{c2}},\qquad X_2=\frac{h_2}{h_{c2}},\qquad h_{c2}=\sqrt[3]{\frac{Q^2}{g\,b_2^2}}
\]
\begin{formula}
$\displaystyle\frac{1}{X_2}-\frac{1}{nX_3}=\frac12\left(X_3^2-X_2^2\right)$\qquad
($X_2$: 1 -- 6;\quad $n$: 1 -- 4.5)
\end{formula}

\section{Opción 2: cepas --- ecuación de Yarnell}
\begin{itemize}
  \item Basada en miles de experimentos realizados entre 1927 y 1936 en
        cauces obstruidos por el efecto de las cepas de diferentes puentes.
  \item Método válido si no se producen resaltos a la salida del puente.
\end{itemize}
\begin{formula}
$\displaystyle\frac{h_3-h_1}{h_1}=\frac{\Delta h}{h_1}=K\cdot Fr_1^2\left(K+5Fr_1^2-0.6\right)\left(\sigma+1.5\sigma^4\right)$
\end{formula}
donde $n$: número de cepas; $b$: ancho del puente [m]; $\sigma=1-\frac{b_e}{b}$:
constante de ancho efectivo (fracción obstruida); $b_e=b-n\cdot D$: ancho
efectivo del puente [m]; $K$: factor de forma de la cepa (sección
transversal), $K=0.9\sim1.15$; $Fr_1$: número de Froude de la sección 1
(aguas abajo).

\begin{figure}[H]
\centering
\begin{tikzpicture}[font=\small,rotate=-20]
  \fill[top color=aguaclara,bottom color=agua] (0,0) rectangle (6,3);
  \fill[gray!40] (2.4,-0.6) rectangle (3.6,3.6);
  \foreach \y in {0.5,1.3,2.1} \fill[gray!70,rounded corners=3pt,draw] (2.3,\y) rectangle (3.7,\y+0.45);
  \foreach \y in {0.35,1.1,1.9,2.75} \draw[azulusm,thick] (0.5,\y) .. controls (2,\y) and (4,\y) .. (5.5,\y);
  \draw[rojo,dashed,thick] (1.8,-0.5) -- (1.8,3.5) node[above]{(3)};
  \draw[rojo,dashed,thick] (4.2,-0.5) -- (4.2,3.5) node[above]{(1)};
  \draw[cota] (5.7,0) -- (5.7,3) node[midway,right]{$b$};
  \draw[cota] (2.1,1.3) -- (2.1,1.75) node[midway,left]{$D$};
\end{tikzpicture}
\caption{Cepas en el cauce (vista en planta).}
\end{figure}

\begin{table}[H]
\centering
\begin{tabular}{lc}
\toprule
\textbf{Forma de la cepa del puente} & \textbf{Coeficiente de Yarnell, $K$}\\
\midrule
Nariz y cola semicircular & 0.90\\
Cepas de doble cilindro con diafragma de conexión & 0.95\\
Cepas de doble cilindro sin diafragma & 1.05\\
Nariz y cola triangular de 90° & 1.05\\
Nariz y cola cuadrada & 1.25\\
Estructura tipo \emph{trestle bent} de diez pilotes & 2.50\\
\bottomrule
\end{tabular}
\caption{Coeficiente de pila de Yarnell.}
\end{table}

\section{Períodos de retorno}
Se utilizan los mismos períodos de retorno del Manual de Carreteras,
Volumen 3, presentados en la clase anterior: para puentes y viaductos,
$T=200$ años (diseño) y 300 años (verificación) en carreteras, y $T=100$ y
150 años en caminos, con vida útil de 50 años.
